\section{Selección de técnicas de aprendizaje de máquina}
La selección de las técnicas de modelado responde a tres criterios complementarios: la consistencia con el alcance metodológico del curso (todas las técnicas seleccionadas están cubiertas en el notebook guía de clase, lo que garantiza coherencia con los conceptos vistos en la asignatura), la pertinencia con el objetivo del negocio definido en la Entrega 1 (identificar qué factores explican el rendimiento académico municipal y descubrir tipologías de municipios para focalizar la inversión pública), y la idoneidad para el tamaño y la naturaleza del dataset (la vista minable Gold cuenta con 7,004 observaciones municipio-año y un conjunto de aproximadamente diez variables predictoras, lo que es manejable para modelos clásicos y se ajusta perfectamente al pipeline de Spark MLlib disponible en el cluster).

\subsection{Técnicas supervisadas: Regresión Lineal y Árbol de Decisión}
Para el modelado supervisado se seleccionaron dos técnicas, ambas presentes en el notebook de clase: \texttt{LinearRegression} y \texttt{DecisionTreeRegressor} del paquete \texttt{pyspark.ml.regression}. La variable objetivo es el puntaje global promedio del municipio (\texttt{avg\_punt\_global}), una variable continua que se busca predecir a partir del perfil contextual del municipio (conectividad, pobreza, cobertura educativa).

La Regresión Lineal se escogió por su capacidad de generar coeficientes interpretables: el modelo entrega directamente el peso marginal de cada variable sobre el puntaje, lo que facilita su uso como insumo para la formulación de políticas públicas. Adicionalmente, al incluir regularización (Ridge mediante \texttt{regParam}, ElasticNet mediante \texttt{elasticNetParam}), el modelo se vuelve robusto frente a multicolinealidad residual entre features. La Regresión Lineal es el modelo base por excelencia y permite establecer un piso de referencia contra el cual evaluar modelos más complejos.

El Árbol de Decisión Regresor se escogió como modelo complementario por su capacidad de capturar interacciones no lineales entre las variables. A diferencia del modelo lineal, el árbol puede modelar relaciones del tipo "si la cobertura es baja Y la conectividad es alta, entonces el efecto sobre el puntaje es distinto al esperado por la suma de los efectos individuales". Esta capacidad es particularmente útil en el contexto del proyecto, donde los hallazgos preliminares (especialmente Q3 y Q8) sugieren que las relaciones entre conectividad, pobreza y rendimiento no son aditivas. Como ventaja adicional, el árbol entrega un ranking de importancia de variables que complementa los coeficientes de la regresión lineal y facilita la interpretación.

La combinación de ambos modelos permite construir una comparativa formal entre el enfoque lineal-paramétrico y el enfoque basado en árboles, ejercicio que se reporta de manera explícita en el notebook \texttt{08\_ml\_supervisado\_regresion} mediante tablas de métricas (MAE, RMSE, $R^2$) sobre el conjunto de test, siguiendo el patrón visto en la Sección 4 del notebook guía.

\subsection{Técnica no supervisada: K-Means}
Para el modelado no supervisado se seleccionó el algoritmo \texttt{KMeans} del paquete \texttt{pyspark.ml.clustering}, también cubierto en el notebook de clase. El objetivo no es predictivo, sino exploratorio: descubrir tipologías de municipios que compartan perfiles similares en términos combinados de conectividad, pobreza y rendimiento.

La elección de K-Means responde a tres razones específicas. Primero, es el algoritmo de clustering más conocido y robusto sobre datos numéricos continuos estandarizados, condiciones que se cumplen en la vista minable después de aplicar \texttt{StandardScaler}. Segundo, los centroides resultantes son directamente interpretables: cada cluster se caracteriza por sus valores promedio en cada feature, lo que permite traducir el resultado matemático en una narrativa de negocio comprensible. Tercero, el algoritmo escala bien sobre Spark y permite probar múltiples valores de $k$ en tiempo razonable, lo que se aprovecha para aplicar el método del codo y el coeficiente de silueta como criterios de selección de la cantidad óptima de clusters.

El uso de K-Means se alinea directamente con el objetivo del negocio: identificar tipologías municipales permite al Ministerio de Educación priorizar las intervenciones de manera focalizada, en lugar de aplicar políticas uniformes que no consideran la heterogeneidad territorial documentada en el análisis exploratorio.

\subsection{Técnicas descartadas y su justificación}
Durante el diseño del proyecto se consideraron otras técnicas que finalmente fueron descartadas para mantener la coherencia con el alcance del curso. Algoritmos como \texttt{RandomForestRegressor} y \texttt{GBTRegressor} (ensembles avanzados), aunque suelen ofrecer mejor desempeño predictivo, no fueron utilizados porque exceden el alcance del notebook guía de clase y no aportarían un beneficio interpretativo proporcional al esfuerzo adicional de explicación en la sustentación oral. De manera similar, se evaluó la posibilidad de incluir modelos de redes neuronales profundas, pero estos se contemplan exclusivamente en la sección de bonos del documento, donde se desarrolla con detalle el modelo \texttt{MultilayerPerceptronClassifier}.

En cuanto a técnicas de clustering alternativas (DBSCAN, Gaussian Mixture Models, clustering jerárquico), se consideraron pero se descartaron porque K-Means es suficiente para el objetivo declarado de tipificación municipal, y porque la presencia de estos algoritmos alternativos en Spark MLlib es menos consistente que la de K-Means.
